%% 01-pendahuluan.tex — bab dokumen contoh skripsi-laporan (hal. 14-15 (proposal), hal. 16-18 (laporan)).
\chapter{PENDAHULUAN}
\section{Latar Belakang}
Kondisi nyata yang ditemui di lapangan belum sesuai dengan kondisi yang diharapkan, sehingga
topik ini penting dan layak diteliti. Uraian diawali dengan fenomena yang ditemui peneliti,
didukung penelitian terdahulu, lalu menunjukkan kesenjangan dan kebaruan penelitian.

Kajian pustaka disusun dari sumber terkini \citep{creswell2018research, patten2017understanding}
dan dari pedoman resmi universitas \citep{unnes2024panduanakademik}. Regulasi penjaminan mutu
menjadi dasar kebijakan \citep{permendikbudristek2023}.
 Contoh rujukan gaya MLA: \citemla[23]{patten2017understanding}.

\begin{table}[htbp]
  \centering
  \caption{Ringkasan penelitian atau kegiatan terdahulu yang relevan}
  \label{tab:terdahulu}
  \begin{tabular}{clc}
    \toprule
    No & Fokus kajian & Hasil utama \\
    \midrule
    1 & Kajian pertama \citep{martin2014write} & Efektif pada ranah kognitif \\
    2 & Kajian kedua \citep{geraldi2017project} & Meningkat kategori sedang \\
    \bottomrule
  \end{tabular}
  \sumber{penulis dari berbagai sumber (2026)}
\end{table}

\section{Rumusan Masalah}
Rumusan masalah dituliskan singkat, padat, dan sistematis dalam bentuk pertanyaan penelitian.

\section{Tujuan Penelitian}
Tujuan penelitian menyatakan upaya penyelesaian masalah yang diuraikan pada rumusan masalah.

\section{Manfaat Penelitian}
Manfaat teoretis berupa sumbangan bagi pengembangan keilmuan, sedangkan manfaat praktis
diarahkan pada pemanfaat hasil penelitian.

\section{Kebaruan Penelitian}
Kebaruan penelitian dijelaskan melalui perbedaan dibandingkan penelitian terdahulu, baik pada
konsep, konteks, metode, maupun penerapannya.
